\documentclass[11pt]{scrreprt}
%%fakesection Load packages
% These are evan.sty
\usepackage{amsmath,amssymb,amsthm}
\usepackage{mathrsfs}
\usepackage[usenames,svgnames,dvipsnames]{xcolor}
\usepackage{hyperref}
\usepackage[nameinlink]{cleveref}
\usepackage{textcomp}
\usepackage{enumerate}
\usepackage[obeyFinal,textsize=scriptsize,shadow]{todonotes}
\usepackage{mathtools}
\usepackage{microtype}

\usepackage[normalem]{ulem}
\usepackage{stmaryrd}
\usepackage{wasysym}
\usepackage{multirow}
\usepackage{prerex}

%%fakesection evan.sty macros
%Small commands
\newcommand{\cbrt}[1]{\sqrt[3]{#1}}
\newcommand{\floor}[1]{\left\lfloor #1 \right\rfloor}
\newcommand{\ceiling}[1]{\left\lceil #1 \right\rceil}
\newcommand{\mailto}[1]{\href{mailto:#1}{\texttt{#1}}}
\newcommand{\ol}{\overline}
\newcommand{\ul}{\underline}
\newcommand{\wt}{\widetilde}
\newcommand{\wh}{\widehat}
\newcommand{\eps}{\varepsilon}
%\renewcommand{\iff}{\Leftrightarrow}
%\renewcommand{\implies}{\Rightarrow}
\newcommand{\vocab}[1]{\textbf{\color{blue} #1}}
\newcommand{\catname}{\mathsf}
\newcommand{\hrulebar}{
  \par\hspace{\fill}\rule{0.95\linewidth}{.7pt}\hspace{\fill}
  \par\nointerlineskip \vspace{\baselineskip}
}
\newcommand{\half}{\frac{1}{2}}

%For use in author command
\newcommand{\plusemail}[1]{\\ \normalfont \texttt{\mailto{#1}}}

%More commands and math operators
\DeclareMathOperator{\cis}{cis}
\DeclareMathOperator*{\lcm}{lcm}
\DeclareMathOperator*{\argmin}{arg min}
\DeclareMathOperator*{\argmax}{arg max}

%Convenient Environments
\newenvironment{soln}{\begin{proof}[Solution]}{\end{proof}}
\newenvironment{parlist}{\begin{inparaenum}[(i)]}{\end{inparaenum}}
\newenvironment{gobble}{\setbox\z@\vbox\bgroup}{\egroup}

%Inequalities
\newcommand{\cycsum}{\sum_{\mathrm{cyc}}}
\newcommand{\symsum}{\sum_{\mathrm{sym}}}
\newcommand{\cycprod}{\prod_{\mathrm{cyc}}}
\newcommand{\symprod}{\prod_{\mathrm{sym}}}

%From H113 "Introduction to Abstract Algebra" at UC Berkeley
\newcommand{\CC}{\mathbb C}
\newcommand{\FF}{\mathbb F}
\newcommand{\NN}{\mathbb N}
\newcommand{\QQ}{\mathbb Q}
\newcommand{\RR}{\mathbb R}
\newcommand{\ZZ}{\mathbb Z}
\newcommand{\charin}{\text{ char }}
\DeclareMathOperator{\sign}{sign}
\DeclareMathOperator{\Aut}{Aut}
\DeclareMathOperator{\Inn}{Inn}
\DeclareMathOperator{\Syl}{Syl}
\DeclareMathOperator{\Gal}{Gal}
\DeclareMathOperator{\GL}{GL} % General linear group
\DeclareMathOperator{\SL}{SL} % Special linear group

%From Kiran Kedlaya's "Geometry Unbound"
\newcommand{\abs}[1]{\left\lvert #1 \right\rvert}
\newcommand{\norm}[1]{\left\lVert #1 \right\rVert}
\newcommand{\dang}{\measuredangle} %% Directed angle
\newcommand{\ray}[1]{\overrightarrow{#1}} 
\newcommand{\seg}[1]{\overline{#1}}
\newcommand{\arc}[1]{\wideparen{#1}}

%From M275 "Topology" at SJSU
\newcommand{\id}{\mathrm{id}}
\newcommand{\taking}[1]{\xrightarrow{#1}}
\newcommand{\inv}{^{-1}}

%From M170 "Introduction to Graph Theory" at SJSU
\DeclareMathOperator{\diam}{diam}
\DeclareMathOperator{\ord}{ord}
\newcommand{\defeq}{\overset{\mathrm{def}}{=}}

%From the USAMO .tex filse
\newcommand{\st}{^{\text{st}}}
\newcommand{\nd}{^{\text{nd}}}
\renewcommand{\th}{^{\text{th}}}
\newcommand{\dg}{^\circ}
\newcommand{\ii}{\item}

% From Math 55 and Math 145 at Harvard
\newenvironment{subproof}[1][Proof]{%
\begin{proof}[#1] \renewcommand{\qedsymbol}{$\blacksquare$}}%
{\end{proof}}

\newcommand{\liff}{\leftrightarrow}
\newcommand{\lthen}{\rightarrow}
\newcommand{\opname}{\operatorname}
\newcommand{\surjto}{\twoheadrightarrow}
\newcommand{\injto}{\hookrightarrow}
\newcommand{\On}{\mathrm{On}} % ordinals
\DeclareMathOperator{\img}{im} % Image
\DeclareMathOperator{\Img}{Im} % Image
\DeclareMathOperator{\coker}{coker} % Cokernel
\DeclareMathOperator{\Coker}{Coker} % Cokernel
\DeclareMathOperator{\Ker}{Ker} % Kernel
\DeclareMathOperator{\rank}{rank}
\DeclareMathOperator{\Spec}{Spec} % spectrum
\DeclareMathOperator{\Tr}{Tr} % trace
\DeclareMathOperator{\pr}{pr} % projection
\DeclareMathOperator{\ext}{ext} % extension
\DeclareMathOperator{\pred}{pred} % predecessor
\DeclareMathOperator{\dom}{dom} % domain
\DeclareMathOperator{\ran}{ran} % range
\DeclareMathOperator{\Hom}{Hom} % homomorphism
\DeclareMathOperator{\End}{End} % endomorphism

% Things Lie
\newcommand{\kb}{\mathfrak b}
\newcommand{\kg}{\mathfrak g}
\newcommand{\kh}{\mathfrak h}
\newcommand{\kn}{\mathfrak n}
\newcommand{\ku}{\mathfrak u}
\newcommand{\kz}{\mathfrak z}
\DeclareMathOperator{\Ext}{Ext} % Ext functor
\DeclareMathOperator{\Tor}{Tor} % Tor functor
\newcommand{\gl}{\opname{\mathfrak{gl}}} % frak gl group
\renewcommand{\sl}{\opname{\mathfrak{sl}}} % frak gl group

% More script letters etc.
\newcommand{\SA}{\mathscr A}
\newcommand{\SB}{\mathscr B}
\newcommand{\SC}{\mathscr C}
\newcommand{\SD}{\mathscr D}
\newcommand{\SF}{\mathscr F}
\newcommand{\SG}{\mathscr G}
\newcommand{\SH}{\mathscr H}
\newcommand{\OO}{\mathcal O}

%% Napkin commands
\newcommand{\prototype}[1]{
	\emph{{\color{red} Prototypical example for this section:} #1} \par\medskip
}
\newenvironment{moral}{%
	\begin{mdframed}[linecolor=green!70!black]%
	\bfseries\color{green!70!black}}%
	{\end{mdframed}}

%%% END EXCERPT OF EVAN.STY
%%fakesection Commutative diagrams
\usepackage{diagrams}
\newarrow{Inj}C--->
\newarrow{Surj}----{>>}
\newarrow{Proj}----{>>}
\newarrow{Embed}>--->
\newarrow{Incl}C--->
\newarrow{Mapsto}|--->
\newarrow{Isom}<--->
\newarrow{Dotted}....>
\newarrow{Dashed}{}{dash}{}{dash}>
\newarrow{Line}-----
\usepackage{tikz-cd}
\usetikzlibrary{decorations.pathmorphing}
\tikzcdset{row sep/normal=3.14em,column sep/normal=3.14em}

%%fakesection Page layout
\usepackage[headsepline]{scrlayer-scrpage}
\renewcommand{\headfont}{}
\addtolength{\textheight}{3.14cm}
\setlength{\footskip}{0.5in}
\setlength{\headsep}{10pt}
\ihead{Evan Chen}
\automark{section}
\chead{}
\ohead{\headmark}
\cfoot{\pagemark}

%%fakesection Fancy section and chapter heads
\addtokomafont{chapterprefix}{\raggedleft}
\renewcommand*{\sectionformat}{\color{purple}\S\thesection\autodot\enskip}

\RedeclareSectionCommand[beforeskip=0.5em]{chapter}
\renewcommand*{\chapterformat}{%
\mbox{\scalebox{1.5}{\chapappifchapterprefix{\nobreakspace}}%
\scalebox{2.718}{\color{purple}\thechapter\autodot}\enskip}}

\addtokomafont{partprefix}{\rmfamily}
\renewcommand*{\partformat}{\color{purple}\scalebox{2.5}{\thepart}%
	\enlargethispage{2em}}
\newcommand{\problemhead}{Problems to think about}

%%fakesection Mini ToC
\usepackage[tight]{minitoc}
\noptcrule
\mtcsetfont{parttoc}{section}{\small\rmfamily\upshape\mdseries}
\renewcommand*{\partheadstartvskip}{}
\renewcommand*{\partheadendvskip}{}
\renewcommand\beforeparttoc{\hspace{\fill}\rule{0.95\linewidth}{2pt}\hspace{\fill}}
\let\oldpart\part
\renewcommand{\part}[1]{\oldpart{#1}\parttoc\cleardoublepage}

%%fakesection SNSD
\newcommand{\snsd}[2][]%
{\includegraphics[#1]{/home/evan/Pictures/TopologicalGG/#2}}

%%fakesection Theorems
\usepackage{thmtools}
\usepackage[framemethod=TikZ]{mdframed}

\theoremstyle{definition}
\mdfdefinestyle{mdbluebox}{%
	roundcorner = 10pt,
	linewidth=1pt,
	skipabove=12pt,
	innerbottommargin=9pt,
	skipbelow=2pt,
	nobreak=true,
	linecolor=blue,
	backgroundcolor=TealBlue!5,
}
\declaretheoremstyle[
	headfont=\sffamily\bfseries\color{MidnightBlue},
	mdframed={style=mdbluebox},
	headpunct={\\[3pt]},
	postheadspace={0pt}
]{thmbluebox}

\mdfdefinestyle{mdrecbox}{%
	linewidth=0.5pt,
	skipabove=12pt,
	frametitleaboveskip=5pt,
	frametitlebelowskip=0pt,
	skipbelow=2pt,
	frametitlefont=\bfseries,
	innertopmargin=4pt,
	innerbottommargin=8pt,
	nobreak=true,
	linecolor=RawSienna,
	backgroundcolor=Salmon!5,
}
\declaretheoremstyle[
	headfont=\bfseries\color{RawSienna},
	mdframed={style=mdrecbox},
	headpunct={\\[3pt]},
	postheadspace={0pt},
]{thmrecbox}

\declaretheorem[%
style=thmbluebox,name=Theorem,numberwithin=section]{theorem}
\declaretheorem[style=thmbluebox,name=Lemma,sibling=theorem]{lemma}
\declaretheorem[style=thmbluebox,name=Proposition,sibling=theorem]{proposition}
\declaretheorem[style=thmbluebox,name=Corollary,sibling=theorem]{corollary}
\declaretheorem[style=thmrecbox,name=Example,sibling=theorem]{example}

\mdfdefinestyle{mdinlinebox}{%
	skipabove=8pt,
	linewidth=3pt,
	rightline=false,
	leftline=true,
	topline=false,
	bottomline=false,
	linecolor=gray,
	backgroundcolor=black!2
}
\declaretheoremstyle[
	headfont=\bfseries,
	bodyfont=\normalfont\small,
	spaceabove=0pt,
	spacebelow=0pt,
	mdframed={style=mdinlinebox}
]{thminlinebox}
\theoremstyle{theorem}
\declaretheorem[name=Question,sibling=theorem,style=thminlinebox]{ques}
\declaretheorem[name=Exercise,sibling=theorem,style=thminlinebox]{exercise}

\theoremstyle{definition}
\newtheorem{claim}[theorem]{Claim}
\newtheorem{definition}[theorem]{Definition}
\newtheorem{fact}[theorem]{Fact}
\newtheorem{remark}[theorem]{Remark}
\newtheorem{abuse}[theorem]{Abuse of Notation}

\newtheorem{problem}{Problem}[chapter]
\renewcommand{\theproblem}{\thechapter\Alph{problem}}
\newtheorem{sproblem}[problem]{Problem}
\newtheorem{dproblem}[problem]{Problem}
\renewcommand{\thesproblem}{\theproblem$^{\star}$}
\renewcommand{\thedproblem}{\theproblem$^{\dagger}$}
\newcommand{\listhack}{$\empty$\vspace{-2em}}

%%fakesection Answers
\usepackage{answers}
\Newassociation{hint}{answeritem}{tex/backmatter/all-hints}
\Newassociation{sol}{answeritem}{tex/backmatter/all-solns}
\renewcommand{\solutionextension}{out}
\renewenvironment{answeritem}[1]{\item[\bf #1.]}{}

%%fakesection Hot chili peppers
\reversemarginpar
\newcommand{\prechili}{\vspace*{0.3em}\hspace*{1.5em}}
\newcommand{\nochili}{\hspace*{1.5em}}
\newcommand{\chili}{\includegraphics[width=1.5em]{media/chili.png}}
\newcommand{\gim}{\marginpar{\prechili\nochili\nochili\chili}}
\newcommand{\yod}{\marginpar{\prechili\nochili\chili\chili}}
\newcommand{\kurumi}{\marginpar{\prechili\chili\chili\chili}}

%%fakesection Table of contents
% First add ToC to ToC
\makeatletter
\usepackage{etoolbox}
\pretocmd{\tableofcontents}{%
  \if@openright\cleardoublepage\else\clearpage\fi
  \pdfbookmark[0]{\contentsname}{toc}%
}{}{}%
\makeatother
\setcounter{tocdepth}{1}
% Fix spacing issues in To
\usepackage[tocindentauto]{tocstyle}
\usetocstyle{KOMAlike}

%%fakesection Asymptote definitions
\usepackage{patch-asy}
\numberwithin{asy}{chapter}
\renewcommand{\theasy}{\thechapter\Alph{asy}}
\begin{asydef}
	import olympiad;
	import cse5;
	size(6cm);
	usepackage("amsmath");
	usepackage("amssymb");
	defaultpen(fontsize(11pt));
	pointfontsize=11;
	void bigbox(string L) {
		filldraw((-4,3)--(4,3)--(4,-3)--(-4,-3)--cycle,
			lightblue+opacity(0.2), black);
		label(L, (4,3), dir(-45));
	}
	void bigblob(string L) {
		filldraw((-4,3)..(-5,1)..(-4.7,0)..(-4.1,-2)..(-4,-3)
			..(-1,-3.3)..(1,-3.2)..(4,-3)
			..(4.3,-1)..(4.1,0.5)..(4,3)
			..(2,3.3)..(0,3.1)..(-3,2.9)..cycle,
			lightblue+opacity(0.2), black);
		label(L, (4,3), dir(5));
	}

	real[] default_ticks = {-4,-3,-2,-1,0,1,2.3,4};
	void cplane(real theta=-45, real[] T = default_ticks, real xmin = -4.5, real xmax = 4.5, real S = 3) {
		graph.xaxis("Re", xmin, xmax, graph.Ticks(format="%", Ticks=T, Size=S), Arrows);
		graph.yaxis("Im", xmin, xmax, graph.Ticks(format="%", Ticks=T, Size=S), Arrows);
		Drawing("0", origin, dir(theta));
		markscalefactor *= xmax;
	}

	pair opendot(pair P, pen p = defaultpen) {
		dot(P, p, UnFill);
		dot(P, p, Draw);
		return P;
	}
\end{asydef}
\def\asydir{asy}

%%fakesection Bibliography
\usepackage[backend=biber,style=alphabetic]{biblatex}
\DeclareLabelalphaTemplate{
	\labelelement{
		\field[final]{shorthand}
		\field{label}
		\field[strwidth=2,strside=left]{labelname}
	}
	\labelelement{
		\field[strwidth=2,strside=right]{year}    
	}
}
\DeclareFieldFormat{labelalpha}{\textbf{\scriptsize #1}}
\addbibresource{references.bib}
\addbibresource{images.bib}

%%fakesection Napkin Macros
\newcommand{\Zc}[1]{\ZZ_{#1}}
\newcommand{\Zm}[1]{\ZZ_{#1}^\times}
\newcommand{\pre}{^{\text{pre}}}
\newcommand{\normalin}{\trianglelefteq}
\newcommand{\triv}{\mathrm{triv}}
\newcommand{\largeotimes}[1]{\mathop{\otimes}\limits_{#1}}
\newcommand{\Mat}{\mathrm{Mat}}
\newcommand{\PGL}[1]{\mathbf{PGL}_{#1}(\CC)}

\DeclareMathOperator{\Stab}{Stab}
\DeclareMathOperator{\FixPt}{FixPt}
\DeclareMathOperator{\refl}{refl}
\DeclareMathOperator{\Fun}{Fun}
\DeclareMathOperator{\Irrep}{Irrep}
\DeclareMathOperator{\Res}{Res}
\DeclareMathOperator{\Reg}{Reg}
\DeclareMathOperator{\Classes}{Classes}
\newcommand{\FunCl}{\Fun_{\mathrm{class}}}
\newcommand{\Homrep}{\Hom_{\mathrm{rep}}}
\newcommand{\ab}{^\text{ab}} % abelianization
\DeclareMathOperator{\ev}{ev}
\DeclareMathOperator{\Wind}{\mathbf I}
\newcommand{\Ctriv}{\CC_{\mathrm{triv}}}
\newcommand{\Csign}{\CC_{\mathrm{sign}}}

% Alg geom macros
\newcommand{\VV}{\mathcal V}
\newcommand{\Vp}{\mathcal V_{\text{pr}}}
\newcommand{\Aff}{\mathbb A}
\newcommand{\II}{\mathcal I_{\text{rad}}}
\newcommand{\RP}{\mathbb{RP}}
\newcommand{\CP}{\mathbb{CP}}
\newcommand{\Cl}{\mathrm{Cl}}
\newcommand{\restrict}[1]{\bgroup\restriction_{#1}\egroup}
\DeclareMathOperator{\Opens}{OpenSets}
\DeclareMathOperator{\Proj}{Proj}
\DeclareMathOperator{\res}{res}
\newcommand{\mm}{\mathfrak m}
\newcommand{\sh}{^\mathrm{sh}}

\renewcommand{\Re}{\opname{Re}}
\renewcommand{\Im}{\opname{Im}}

%% Category Theory macros
\renewcommand{\AA}{\mathcal A}
\newcommand{\BB}{\mathcal B}
\newcommand{\obj}{\operatorname{obj}}
\newcommand{\op}{^{\mathrm{op}}}


\usepackage{tkz-graph}
\pgfarrowsdeclare{biggertip}{biggertip}{%
  \setlength{\arrowsize}{1pt}  
  \addtolength{\arrowsize}{.1\pgflinewidth}  
  \pgfarrowsrightextend{0}  
  \pgfarrowsleftextend{-5\arrowsize}  
}{%
  \setlength{\arrowsize}{1pt}  
  \addtolength{\arrowsize}{.1\pgflinewidth}  
  \pgfpathmoveto{\pgfpoint{-5\arrowsize}{4\arrowsize}}  
  \pgfpathlineto{\pgfpointorigin}  
  \pgfpathlineto{\pgfpoint{-5\arrowsize}{-4\arrowsize}}  
  \pgfusepathqstroke  
}  
\tikzset{
	EdgeStyle/.style = {>=biggertip}
}

\usepackage[all,cmtip,2cell]{xy}
\UseTwocells
\newcommand\nattfm[5]{\xymatrix@C+2pc{#1 \rtwocell<4>^{#2}_{#4}{\; #3} & #5}}

%% Alg Top macros
\DeclareMathOperator{\Cells}{Cells}
\newcommand{\HdR}{H_{\mathrm{dR}}}

%% Alg NT macros
\renewcommand{\aa}{\mathfrak a}
\newcommand{\bb}{\mathfrak b}
\newcommand{\pp}{\mathfrak p}
\newcommand{\qq}{\mathfrak q}

\newcommand{\Frob}{\mathrm{Frob}}
\DeclareMathOperator{\Norm}{N}
\DeclareMathOperator{\Ram}{Ram}
\newcommand{\TrK}{\Tr_{K/\QQ}}
\newcommand{\NK}{\Norm_{K/\QQ}}
\newcommand{\kf}{\mathfrak f}
\newcommand{\kP}{\mathfrak P}
\newcommand{\kQ}{\mathfrak Q}

%% Diff geo macros
\newcommand{\ee}{\mathbf e}
\newcommand{\fpartial}[2]{\frac{\partial #1}{\partial #2}}

%% ST macros
\newcommand{\CH}{\mathsf{CH}}
\newcommand{\ZFC}{\mathsf{ZFC}}

\newcommand{\Name}{\text{Name}}
\newcommand{\Po}{\mathbb P}
\newcommand{\nrank}{\opname{n-rank}} % ranks
\newcommand{\PP}{\mathcal P}

\newcommand{\EmptySet}{\mathrm{EmptySet}}
\newcommand{\PowerSet}{\mathrm{PowerSet}}
\newcommand{\Pairing}{\mathrm{Pairing}}
\newcommand{\Infinity}{\mathrm{Infinity}}
\newcommand{\Extensionality}{\mathrm{Extensionality}}
\newcommand{\Foundation}{\mathrm{Foundation}}
\newcommand{\Union}{\mathrm{Union}}
\newcommand{\Comprehension}{\mathrm{Comprehension}}
\newcommand{\Replacement}{\mathrm{Replacement}}

\usepackage{mathrsfs}

\newcommand\MM{\mathscr M}
\newcommand\llex{<_{\text{lex}}}

\DeclareMathOperator{\cof}{cof}

%% Quantum macros
\usepackage{braket}
\newcommand{\cvec}[1]{\begin{bmatrix} #1 \end{bmatrix}}
\newcommand{\pair}[2]{\begin{bmatrix} #1 \\ #2 \end{bmatrix}}
\newcommand{\zup}{\ket\uparrow}
\newcommand{\zdown}{\ket\downarrow}
\newcommand{\xup}{\ket\rightarrow}
\newcommand{\xdown}{\ket\leftarrow}
\newcommand{\yup}{\ket\otimes}
\newcommand{\ydown}{\ket\odot}
\newcommand{\UCNOT}{U_{\mathrm{CNOT}}}
\newcommand{\UQFT}{U_{\mathrm{QFT}}}

\def\asydir{}
\addbibresource{../../references.bib}
\renewcommand{\gim}{*}
\renewcommand{\yod}{**}
\renewcommand{\kurumi}{***}

\begin{document}
\title{Algebraic Geometry II: Schemes}
\maketitle

\tableofcontents

\chapter{Sheaves and ringed spaces}
Most of the complexity of the affine variety $V$ earlier comes from $\OO_V$.
This is a type of object called a ``sheaf''.
The purpose of this chapter is to completely define what this sheaf is,
and just what it is doing.

The typical example to keep in mind is a sheaf of
``functions with property $P$'' on a topological space $X$:
for every open set $U$, $\SF(U)$ gives us the ring of functions on $X$.
However, we will work very abstractly and only assume $\SF(U)$
is a ring, without an interpretation as ``functions''.

The payoff for this abstraction is that it will
allow us to define an arbitrary scheme in the next chapter.
Varieties use $\CC[x_1, x_2, \dots, x_n] / I$ as their ``ring of functions'',
and by using the fully general sheaf we will be replace this
with \emph{any} commutative ring.
In particular, we can use the case where $I$ is not radical, such as
$\CC[x] / (x^2)$; this gives the ``multiplicity''
behavior that we sought after all along.  

\section{Pre-sheaves}
\prototype{The sheaf of holomorphic (or regular, continuous,
differentiable, constant, whatever) functions.}

The proper generalization of our $\OO_V$ is a so-called sheaf of rings.
Recall that $\OO_V$ was a took \emph{open sets of $V$} to \emph{rings},
with the interpretation that $\OO_V(U)$ was a ``ring of functions''.

In light of this, we first make the following definition
for the general construct, in the language of categories.

\begin{definition}
	For a topological space $X$ let $\Opens(X)$ denote
	its open sets of $X$.
\end{definition}
\begin{definition}
	A \vocab{pre-sheaf} of rings on a space $X$ is a function
	\[ \SF : \Opens(X) \to \catname{Rings} \]
	meaning each open set gets associated with a ring $\SF(U)$.
	Each individual element of $\SF(U)$ is called a \vocab{section}.
	An element of $\SF(X)$ is called a \vocab{global section}.

	It is also equipped with the following information:
	for any $U_1 \subseteq U_2$ there is a \vocab{restriction map}
	\[ \res_{U_1,U_2} : \SF(U_2) \to \SF(U_1) \]
	such that $\res_{U,U}$ is the identity and
	whenever $U_1 \subseteq U_2 \subseteq U_3$ the diagram
	\begin{diagram}
		\SF(U_3) & \rTo^{\res_{U_2,U_3}} & \SF(U_2) \\
		& \rdTo_{\res_{U_1,U_3}} & \dTo_{\res_{U_1,U_2}} \\
		&& \SF(U_1)
	\end{diagram}
	commutes. (Restricting ``big to medium to small''
	is the same as ``big to small''.)
\end{definition}

\begin{abuse}
	If $s \in \mathscr F(U_2)$ is some section and $U_1 \subseteq U_2$,
	then rather than write $\res_{U_1,U_2}(s)$
	I will write $s\restrict{U_1}$ instead:
	``$s$ restricted to $U_1$''.
	This is abuse of notation because the section $s$ is just
	an element of some ring, and in the most abstract of cases
	may not have a natural interpretation as function.
\end{abuse}

\begin{example}[Examples of pre-sheaves]
	\listhack
	\begin{enumerate}[(a)]
		\ii For an affine variety $V$, $\OO_V$ is of course a sheaf,
		with $\OO_V(U)$ being the ring of regular functions on $U$.
		The restriction map just says that if $U_1 \subseteq U_2$,
		then a function $s \in \OO_V(U_2)$ can also be thought of as
		a function $s \restrict{U_1} \in \OO_V(U_1)$,
		hence the name ``restriction''.
		The commutativity of the diagram then follows.
		
		\ii Let $X \subseteq \RR^n$ be an open set.
		Then there is a sheaf of smooth/differentiable/etc.\ functions on $X$.
		In fact, one can do the same construction for any manifold $M$.

		\ii Similarly, if $X \subseteq \CC$ is open,
		we can construct a sheaf of holomorphic functions on $X$.
	\end{enumerate}
	In all these examples, the sections $s \in \SF(U)$
	are really functions on the space, but in general they need not be.
\end{example}

Now, we give a second, equivalent and far shorter definition of pre-sheaf:
\begin{abuse}
	By abuse of notation, $\Opens(X)$ will also be thought of as a
	posetal category by inclusion. Thus $\varnothing$ is an initial object
	and the entire space $X$ is a terminal object.
\end{abuse}
\begin{definition}
	A \vocab{pre-sheaf} of rings on $X$ is a contravariant functor
	\[ \SF : \Opens(X)\op \to \catname{Rings}. \]
\end{definition}
\begin{ques}
	Check that these definitions are equivalent.
\end{ques}
It is now clear that we can actually replace $\catname{Rings}$
with any category we want.

So we should think:
\begin{moral}
	Pre-sheaves should be thought of as
	``returning the ring of functions with a property $P$''.
\end{moral}

\section{Sheaves}
\prototype{Constant functions aren't pre-sheaves,
	but locally constant ones are.}

Now, the main idea next is that
\begin{moral}
	Sheaves are pre-sheaves for which $P$ is a \emph{local} property.
\end{moral}

The formal definition doesn't illuminate this as much as the examples
do, but I have to give it first for the examples to make sense.

\begin{definition}
	A \vocab{sheaf} $\mathscr F$ is a pre-sheaf satisfying the following
	two additional axioms:
	Suppose $U$ is covered by open sets $U_\alpha \subseteq U$. Then:
	\begin{enumerate}
		\ii (Identity) If $s, t \in \mathscr F(U)$ are sections,
		and $s\restrict{U_\alpha} = t\restrict{U_\alpha}$
		for all $\alpha$, then $s = t$.
		\ii (Collation) Consider sections
		$s_\alpha \in \mathscr(U_\alpha)$ for each $\alpha$.
		Suppose that 
		\[ s_\alpha \restrict{U_\alpha \cap U_\beta}
			= s_\beta \restrict{U_\alpha \cap U_\beta} \]
		for each $U_\alpha$ and $U_\beta$.
		Then we can find $s \in U$ such that
		$s \restrict{U_\alpha}  = s_\alpha$.
	\end{enumerate}
\end{definition}
This is best illustrated by picture: consider an open cover $U_1 \cup U_2$.
\begin{center}
	\begin{asy}
		size(4cm);
		filldraw(shift(-0.5,0)*unitcircle, lightred+opacity(0.3), red);
		filldraw(shift( 0.5,0)*unitcircle, lightblue+opacity(0.3), blue);
		label("$U_1$", (-0.5,0)+dir(135), dir(135), red);
		label("$U_2$", ( 0.5,0)+dir(45), dir(45), blue);
	\end{asy}
\end{center}
Then for a sheaf of functions, the axioms are saying that:
\begin{itemize}
	\ii If $s$ and $t$ are functions (with property $P$)
	on the whole space $U = U_1 \cup U_2$,
	and $s \restrict{U_1} = t \restrict{U_1}$,
	$s \restrict{U_2} = t \restrict{U_2}$,
	then $s = t$ on the entire union.
	This is clear.

	\ii If $s_1$ is a function with property $P$ on $U_1$
	and $s_2$ is a function with property $P$ on $U_2$,
	and the two functions agree on the overlap,
	then one can collate them to obtain a function $s$
	on the whole space:
	this is obvious, but \textbf{the catch is that the collated function
	needs to have property $P$ as well.}
	That's why it matters that $P$ is local.
\end{itemize}

The reason we need these axioms is that in our abstract definition of a sheaf,
the output of the sheaf is an abstract ring, which need not actually
have a concrete interpretation as ``functions on $X$'', even though
our examples will usually have this property.

Now for the examples, which are more enlightening:
\begin{example}
	[Examples and non-examples of sheaves]
	\listhack
	\begin{enumerate}[(a)]
		\ii Pre-sheaves of arbitrary / continuous / differentiable / smooth
		/ holomorphic functions are still sheaves.
		This is because to verify a function is continuous,
		one only needs to look at small neighborhoods at once.
		
		\ii For a complex variety $V$, $\OO_V$ is a sheaf,
		precisely because our definition was \emph{locally} quotients
		of polynomials.

		\ii The pre-sheaf of \emph{constant} real functions on a space $X$
		is \emph{not} a sheaf, because it fails the collation axiom.
		Namely, suppose that $U_1 \cap U_2 = \varnothing$.
		Then if $s_1$ is the constant function $1$ on $U_1$
		while $s_2$ is the constant function $2$ on $U_2$,
		then we cannot collate these to a constant function on $U_1 \cup U_2$.

		\ii On the other hand, \emph{locally constant} functions
		do produce a sheaf. (A function is locally constant
		if for every point it is constant on some neighborhood.)
	\end{enumerate}
	In fact, the sheaf in (c) is what is called a \emph{sheafification}
	of the pre-sheaf constant functions, which we define momentarily.
\end{example}

\section{Stalks}
\prototype{Germs of real smooth functions tell you the derivatives,
but germs of holomorphic functions determine the entire function.}
Let $\SF$ be a pre-sheaf.
If we have a function $s \in \SF(U)$ and a point $p \in U$,
then in general it doesn't make sense to ask what $s(p)$ is
(even though all our examples look like this),
because $\SF(U)$ is an arbitrary ring.
So, we will replace the notion of $s(p)$ with a so-called \emph{germ}.

\begin{definition}
	Let $\SF$ be a pre-sheaf of rings.
	For every point $p$ we define the \vocab{stalk} $\SF_p$ to be the set
	\[ \left\{ (s, U) \mid s \in \SF(U), p \in U \right\} \]
	modulo the relation $\sim$ that 
	\[ (s_1,U_1) \sim (s_2, U_2) \text{ if }
		s_1 \restrict{U_1 \cap U_2} = s_2 \restrict{U_1 \cap U_2}. \]
	The equivalence classes themselves are called \vocab{germs}.
\end{definition}

\begin{definition}
	The germ of a given $s \in \SF(U)$ at a point $p$
	is the equivalence class for $(s,U) \in \SF_p$.
	We denote this by $[s]_p$.
\end{definition}

So what's happening is: We consider functions $s$ defined near $p$
and, since we only care about local behavior,
we identify any two functions agreeing on a neighborhood of $p$,
no matter how small.

Notice that the stalk is itself a ring as well:
for example, addition is done by
\[ 
	\left( s_1, U_1 \right) + \left( s_2, U_2 \right)
	=
	\left( s_1 \restrict{U_1 \cap U_2} + s_2 \restrict{U_1 \cap U_2},
	U_1 \cap U_2 \right).
\]

So the germ $[s]_p$ now plays the role of the ``value'' $s(p)$.
But actually, it carries much more information than that.

\begin{example}
	[Germs of real smooth functions]
	Let $X = \RR$ and let $\SF$ be the sheaf on $X$ of smooth functions
	(i.e.\ $\SF(U)$ is the set of smooth real-valued functions on $U$).

	Consider a global section, $s : \RR \to \RR$ (thus $s \in \SF(X)$)
	and its germ at $0$.
	\begin{enumerate}[(a)]
		\ii From the germ we can read off $s(0)$, obviously.
		\ii We can also find $s'(0)$, because the germ carries enough
		information to compute the limit $\lim_{h \to 0} \frac1h[s(h)-s(0)]$.
		\ii Similarly, we can compute the second derivative and so on.
		\ii However, we can't read off, say, $s(3)$ from the germ.
		For example, take
		\[
			s(x) = \begin{cases}
				e^{-\frac{1}{x-1}} & x > 1 \\
				0 & x \le 1.
			\end{cases}
		\]
		Note $s(3) = e^{-\half}$, but $[\text{zero function}]_0 = [s]_0$.
		So germs can't distinguish between the zero function and $s$.
	\end{enumerate}
\end{example}
\begin{example}
	[Germs of holomorphic functions]
	Holomorphic functions are very strange in this respect.
	Consider the sheaf $\SF$ on $\CC$ of \emph{holomorphic} functions.

	Take $s : \CC \to \CC$ a global section.
	Given the germ of $s$ at $0$, we can read off $s(0)$, $s'(0)$, et cetera.
	The miracle of complex analysis is that just knowing
	the derivatives of $s$ at zero is enough to reconstruct all of $s$:
	we can compute the Taylor series of $s$ now.
	\textbf{Thus germs of holomorphic functions determine the entire function};
	they ``carry much more information'' than their real counterparts.
	
	In particular, we can concretely describe the sheaf:
	\[
		\SF_p = \left\{
			\sum_{k \ge 0} c_k (z-p)^k
			\text{ convergent near $p$}
		\right\}.
	\]
	In particular, this includes germs of meromorphic functions,
	so long as there is no pole at $p$ itself.
\end{example}

And of course, our algebraic geometry example.
\begin{abuse}
	Rather than writing $(\OO_V)_p$ we will write $\OO_{V,p}$.
\end{abuse}
\begin{example}
	[Stalks of the sheaf on an affine variety]
	Let $V \subseteq \Aff^n$ be a variety, and assume $p \in V$.
	Then, a regular function $\varphi$ on $U \subseteq V$
	is supposed to be a function on $U$ that ``locally'' is a quotient
	of two functions in $\CC[V]$.

	However, \textbf{as far as the germ is concerned, we only care about
	whichever quotient applies near the point $p$}.
	In light of this, we only think about the representation at $p$,
	and ignore the local clause completely:
	thus the entire stalk can be thought of as
	\[
		\OO_{V,p} = 
		\left\{ \left( \tfrac fg , U \right) \mid 
			U \ni p, \; f,g \in \CC[V], \;
			\text{$g \neq 0$ on $U$} \right\}
	\]
	modulo the usual relations.

	Now, since we happen to be working with complex polynomials,
	we know that a rational function is determined by its
	behavior on any neighborhood of $p$
	(complex analysis forever!).
	Thus 
	\[
		\OO_{V,p} =
		\left\{ \tfrac fg \mid f,g \in \CC[V], \; g(p) \neq 0 \right\}.
	\]
	which don't vanish on $p$.
\end{example}

\begin{remark}
	[For category lovers]
	You may notice that $\SF_p$ seems to be
	``all the $\SF_p(U)$ coming together'', where $p \in U$.
	And in fact, $\SF_p(U)$ is the categorical \emph{limit}
	of the diagram formed by all the $\SF(U)$ such that $p \in U$.
	This is often written
	\[ \SF_p = \varinjlim_{U \ni p} \SF(U) \]
	Thus we can define stalks in any category with limits,
	though to be able to talk about germs the category needs
	to be concrete.
\end{remark}

\section{Sections ``are'' sequences of germs}
\prototype{A real function on $U$ is a sequence of
	real numbers $f(p)$ for each $p \in U$ satisfying some local condition.
	Analogously, a section $s \in \SF(U)$ is a sequence of germs
	satisfying some local compatibility condiiton.
}
Let $\SF$ be a sheaf on a space $X$ now.
The purpose of this section is to convince you that the correct way 
think about a section $s \in \SF(U)$ is as a sequence of germs
above every point $p \in U$.

How do we do this?
Given $s \in \SF(U)$ we consider its germ $[s]_p$ above every for $p \in U$.
To visualize, picture an open set $U \subseteq X$,
and above every point $p \in P$ imagine there is a little speck $[s]_p$,
which bundles the information of $s$ near $p$.
\begin{example}[Real functions vs.\ germs]
	Let $X$ be a space and let $\SF$ be the sheaf of smooth functions 
	Given a section $f \in \SF(U)$,
	\begin{itemize}
		\ii As a function, $f$ is just a choice of value $f(p) \in \RR$ at
		every point $p$, subject to a local ``smooth'' condition.

		\ii Let's now think of $f$ as a sequence of germs.
		At every point $p$ the germ $[f]_p \in \SF_p$ gives us the value $f(p)$
		as we described above. The germ packages even more data than this:
		from the germ $[f]_p$ alone we can for example compute $f'(p)$.
		Nonetheless we stretch the analogy and think of $f$
		as a choice of germ $[f]_p \in \SF_p$ at each point $p$.
	\end{itemize}
	Thus we can replace the notion of the value $f(p)$ with germ $[f]_p$.
	This is useful because in a general sheaf $\SF$, the notion $s(p)$
	is not defined while the notion $[s]_p$ is.
\end{example}

From the above example it's obvious that if we know each germ $[s]_p$
this should let us reconstruct the entire section $s$.
Let's check this from the sheaf axioms:
\begin{exercise}
	[Sections are determined by stalks]
	Let $\SF$ be a sheaf.
	Consider the natural map $\SF(U) \to \prod_{p \in U} \SF_p$ described above.
	Show that this map is injective, i.e.\
	the germs of $s$ at every point $p \in U$ determine the section $s$.
	(You will need the ``identity'' sheaf axiom, but not ``collating''.)
\end{exercise}

However, this map is clearly not surjective!
\begin{ques}
	Come up with a counterexample to break surjectivity.
	(This is like asking ``come up with a non-smooth function''.)
\end{ques}
Nonetheless we can describe the image:
we want a sequence of germs $(g_p)_{p \in U}$ such that near every germ $g_p$,
the germs $g_q$ are ``compatible'' with $g_p$.
We make this precise:
\begin{definition}
	Let $\SF$ be sheaf (or even a pre-sheaf) and $U$ open.
	A sequence $(g_p)_{p \in U}$ of germs
	($g_p \in \SF_p$) is said to be \vocab{compatible} if
	they can be ``locally collated'':
	\begin{quote}
		For any $p \in U$ there exists a neighborhood $U_p \ni p$
		and a section $s \in \SF(U_p)$ on it
		such that $[s]_q = g_q$ for each $q \in U_p$.
	\end{quote}
	Intuitively, the germs should ``collate together'' to some section near
	each \emph{individual} point $q$
	(but not necessarily to a section on all of $U$).
\end{definition}

We let the reader check this definition is what we want:
\begin{exercise}
	Prove that any choice of compatible germs over $U$
	collates together to a section of $U$.
	(You will need the ``collating'' sheaf axiom, but not ``identity''.)
\end{exercise}

Putting together the previous two exercise gives:
\begin{theorem}
	[Sections ``are'' just compatible germs]
	Let $\SF$ be a sheaf.
	There is a natural bijection between
	\begin{itemize}
		\ii sections of $\SF(U)$, and
		\ii sequences of compatible germs over $U$.
	\end{itemize}
\end{theorem}
This is in exact analogy to the way that e.g.\
a smooth real-valued function on $U$ is a choice
of real number $f(p) \in \RR$ at each point $p \in U$
satisfying a local smoothness condition.

Thus the notion of stalks is what lets us recover the viewpoint
that sections are ``functions''.  Therefore for theoretical purposes,
\begin{moral}
	You should usually think of a section $s \in \SF(U)$ as a sequence of germs.
\end{moral}
In particular, this makes restriction morphisms easy to deal with:
just truncate the sequence of germs!

\section{Sheafification}
\prototype{The pre-sheaf of locally constant functions
	becomes the sheaf of constant functions.}

Now that we have the language of germs,
we can define the so-called sheafification.
The idea is that if $\SF$ is the pre-sheaf of ``functions with property $P$''
then we want to associate a sheaf $\SF\sh$ of
``functions which are locally $P$'', which makes them into a sheaf.
We have already seen two examples of this:
\begin{example}
	[Sheafification]
	\listhack
	\begin{enumerate}[(a)]
		\ii If $X$ is a topological space,
		and $\SF$ is the pre-sheaf of constant functions on open sets of $X$,
		then $\SF\sh$ is the sheaf of locally constant functions.

		\ii If $V$ is an affine variety,
		and $\SF$ is the pre-sheaf of rational functions,
		then $\SF\sh$ is the sheaf of regular functions
		(which are locally rational).
	\end{enumerate}
\end{example}

So how do we encode ``locally $P$''?
Answer: we saw last time that for a sheaf $\SF$,
compatible germs biject to sections.
So we \emph{force} this to be true by defining:
\begin{definition}
	The \vocab{sheafification} $\SF\sh$ of a pre-sheaf $\SF$ is defined by
	\[ \SF\sh(U) =
		\left\{ \text{sequences of compatible
		germs $(g_p)_{p \in U}$} \right\}.  \]
\end{definition}
\begin{exercise}
	Reconcile this definition with the two examples we gave.
\end{exercise}
\begin{abuse}
	Technically I haven't told you what the restriction morphisms
	of $\SF\sh(U)$ are but, they are not hard to guess.
	I'll usually be equally sloppy in the future:
	when defining a sheaf $\SF$, I'll only say what $\SF(U)$ is,
	with the restriction morphisms $\SF(U_2) \to \SF(U_1)$ being implicit.
\end{abuse}
The construction is contrived so that given a section
$(g_p)_{p \in U} \in \SF\sh(U)$ the germ at a point $p$ is $g_p$:
\begin{lemma}
	[Pre-sheaves and sheaves have the same stalk]
	\label{lem:pre_sheaf_stalk}
	Let $\SF$ be a pre-sheaf and $\SF\sh$ its sheafification.
	Let $q \in U$ with $U$ an open set.
	Then there is an isomorphism
	\[ (\SF\sh)_q \cong \SF_q. \]
\end{lemma}
\begin{proof}
	A germ in $(\SF\sh)_q$ looks like 
	$\left( (g_p)_{p \in U}, U \right)$,
	where $g_p = (s_p, U_p)$ are themselves germs of $\SF_p$.
	Then the isomorphism is given by
	\[ \left( (g_p)_{p \in U}, U \right) \mapsto g_q \in \SF_q. \]
	The inverse map is given by for each $g = (s,U) \in \SF_q$ by
	\[ g \mapsto \left( (g)_{p \in U}, U \right) \in (\SF\sh)_q \]
	i.e.\ the sequence of germs is the constant sequence.
\end{proof}

We'll later see that if $\SF$ is already a sheaf, then $\SF\sh = \SF$.

\section{Morphisms of sheaves}
First, recall that a sheaf is a contravariant functor (pre-sheaf)
with extra conditions. In light of this, it is not hard to guess
the definition of a morphism of pre-sheaves:
\begin{definition}
	A \vocab{morphism of (pre-)sheaves} $\alpha : \SF \to \SG$ on the same
	space $X$ is a \textbf{natural transformation} of the underlying functors.
	Isomorphism of sheaves is defined in the usual way.
\end{definition}
\begin{ques}
	Show that this amounts to: for each $U \subseteq X$ we need to specify
	a morphism $\alpha_U : \SF(U) \to \SG(U)$ such that the diagram
	\begin{diagram}
		\SF(U_2) & \rTo^{\alpha_{U_2}} & \SG(U_2) \\
		\dTo^{\res_{U_1, U_2}} && \dTo_{\res_{U_1, U_2}} \\
		\SF(U_1) & \rTo_{\alpha_{U_1}} & \SG(U_1)
	\end{diagram}
	commutes any time that $U_1 \subseteq U_2$.
\end{ques}
However, in the sheaf case we like stalks more than sections because
they are theoretically easier to think about.
And in fact:
\begin{proposition}
	[Morphisms determined by stalks]
	A morphism of sheaves $\alpha : \SF \to \SG$ induces a morphism of stalks
	\[ \alpha_p : \SF_p \to \SG_p \]
	for every point $p \in X$.
	Moreover, the sequence $(\alpha_p)_{p \in X}$ determines $\alpha$ uniquely.
\end{proposition}
\begin{proof}
	The morphism $\alpha_p$ itself is just
	$(s, U) \xmapsto{\alpha_p} (\alpha_U(s), U)$.
	\begin{ques}
		Show this is well-defined.
	\end{ques}
	Now suppose $\alpha , \beta : \SF \to \SG$ satisfy $\alpha_p = \beta_p$
	for every $p$. We want to show $\alpha_U(s) = \beta_U(s)$
	for every $s \in U$.
	\begin{ques}
		Verify this using the description of sections
		as sequences of germs. \qedhere
	\end{ques}
\end{proof}
Thus a morphism of sheaves can be instead modelled as a morphism
of all the stalks. We will see later on that this viewpoint is quite useful.


\section{Local rings, and locally ringed spaces}
\prototype{Smooth functions $f$ on $X \subseteq \RR^n$ have invertible
germs at $p$ unless $f(p) = 0$.}
The stalks of the examples we produced above are special types
of rings, called \emph{local rings}.
Algebraically, the definition of these is:
\begin{definition}
	A \vocab{local ring} $R$ is a ring with exactly one maximal ideal.
\end{definition}
\begin{exercise}
	Show that a ring $R$ is a local ring if there exists
	a proper ideal $\mm \subsetneq R$ such that
	all elements of $R \setminus \mm$ are units.
	(Bonus: prove the converse.)
\end{exercise}

% Wikipedia has a good explanation here
To see why this definition applies to the stalks above,
we need to identify what the maximal ideal is.
Let's go back to the example of $X = \RR$ and $\SF(U)$ the smooth functions,
and consider the stalk $\SF_{p}$, where $p \in X$.
Define the ideal $\mm_p$ to be the set of germs $(s,U)$ for which $s(p) = 0$.

Then $\mm_p$ is maximal: we have an exact sequence
\[ 0 \to \mm_p \to \SF(U) \taking{(s,U) \mapsto s(p)} \RR \to 0 \]
and so $\SF(U) / \mm_p \cong \RR$, which is a field.

It remains to check there are no nonzero maximal ideals.
Now note that if $s \notin \mm_p$,
then $s$ is nonzero in some neighborhood of $p$,
then one can construct the function $1/s$ in a neighborhood of $p$.
So \textbf{every element of $\SF_p \setminus \mm_p$ is a unit};
$\mm_p$ is in fact the only maximal ideal!

More generally,
\begin{moral}
	If $\SF$ consists of ``field-valued functions'',
	the stalk $\SF_p$ probably has a maximal ideal
	consisting of the germs vanishing at $p$.
\end{moral}
The discussion above implies, for example, the following.
\begin{proposition}
	[Stalks are often local rings]
	The stalks of each of the following types of sheaves are local rings:
	\begin{enumerate}[(a)]
		\ii Sheaves of continuous real/complex functions on a topological space
		\ii Sheaves of smooth functions on any manifold
		\ii Regular functions on an algebraic variety $V$.
	\end{enumerate}
\end{proposition}

We can now define:
\begin{definition}
	A \vocab{ringed space} is a topological space $X$ equipped
	with a sheaf $\OO_X$ of rings.
	Suppose that for every point $p$, the stalk $\OO_{X,p}$
	is a local ring.
	Then we say that $\OO_X$ is a \vocab{locally ringed space}.
	We denote the maximal ideals by $\mm_{X,p}$.
\end{definition}

In particular, in the previous chapter we showed that every
affine variety could be built into a locally ringed space. Hooray!
\begin{abuse}
	A ringed space $(X, \OO_X)$ is abbreviated to just $X$,
	while $p \in X$ means ``$p$ is in the topological space $X$''.
\end{abuse}


\section{(Optional) Defining the morphisms}
Finally, it remains to define a morphism of locally ringed space.
To do this we have to build up in several steps.

\begin{remark}
	This is marked as optional because it involves a \emph{lot}
	of definitions, and because we will more or less
	ignore the actual definition later on when we start applying
	schemes to things. One can probably get away with just
	suspending belief and knowing ``there is a reasonable definition
	of morphisms of locally ringed spaces''.
\end{remark}

\subsection*{Morphisms of ringed spaces}
Suppose we have ringed spaces $X = (X, \OO_X)$ and $Y = (Y, \OO_Y)$,
and we want to define a map
\[ f : X \to Y. \]

On the level of spaces, $f : X \to Y$ should of course be a continuous map.
But we also want to $f$ to do something with $\OO_X$ and $\OO_Y$.
However, we only defined how to take a morphism of sheaves
that act on the \emph{same} space!
Thus we impose the following definition,
which lets us \emph{push} the sheaf on $X$ to a sheaf on $Y$.
Picture:
\begin{diagram}
	X && \rTo^f & Y && X && \rTo^f & Y \\
	\Opens(X)\op && \lTo^{f\pre} & \Opens(Y)\op
	&& f\pre(U) && \lMapsto^{f\pre} & U & \\
	& \rdMapsto(3,2)_{\SF} && \dMapsto_{f_\ast \SF}
		&& & \rdMapsto(3,2)_{\SF} && \dMapsto_{f_\ast \SF} & \\
	&&& \catname{Rings}
		&& && \SF(f\pre(U)) & = (f_\ast \SF)(U) & \in \catname{Rings}
\end{diagram}

\begin{definition}
	Let $\SF$ be a sheaf on $X$, and $f : X \to Y$ a continuous map.
	The \vocab{pushforward sheaf} $f_\ast \SF$ on $Y$ is defined by
	\[ (f_\ast \SF)(U) = \SF(f\pre(U)) \qquad \forall U \subseteq Y. \]
	This makes sense, since $f\pre(U)$ is open in $X$.
\end{definition}
As $f\pre$ is a functor $\Opens(Y)\op \to \Opens(X)\op$,
\textbf{the pushforward $f_\ast\SF$ is
just the composition of these two functors}.

\begin{ques}
	Technically $f_\ast\SF$ is supposed to be a
	functor $\Opens(Y)\op \to \catname{Rings}$, so it also needs
	to come with restriction arrows. What are they?
\end{ques}

I haven't actually checked that $f_\ast \SF$ is a sheaf
(as opposed to a pre-sheaf), but this isn't hard to do.

Now we can define a morphism of ringed spaces.
In addition to the topological spaces,
we also include a natural transformation between
the two structure sheafs, after pushing the one on $X$ forward.
\begin{definition}
	A \vocab{morphism of ringed spaces} $f : (X, \OO_X) \to (Y, \OO_Y)$
	consists of a continuous map of topological spaces $f : X \to Y$,
	as well as a morphism of sheaves $f^\ast : \OO_Y \to (f_\ast \OO_X)$.
\end{definition}
The latter is thus a map
$f^\ast_U : \OO_Y(U) \to (f_\ast \OO_X)(U) = \OO_X(f\pre(U))$
for every $U \subseteq Y$, satisfying
the axioms for a natural transformation.

\subsection*{Morphisms of locally ringed spaces}
Last step!
Suppose now that $(X, \OO_X)$ and $(Y, \OO_Y)$ are locally ringed spaces.
Thus we need to deal with some information about the stalks.

Given a morphism of ringed spaces $f : (X, \OO_X) \to (Y, \OO_Y)$,
we can actually use the $f^\ast_U$ above to induce maps on the stalks,
as follows. For a point $p \in X$, construct the map
\begin{diagram}
	f^\ast_P :& \OO_{Y, f(p)} & \rTo && \OO_{X, f(p)} \\
	& (s, U) & \rMapsto && (f^\ast_U(s), f\pre(U)).
\end{diagram}
where $s \in \OO_Y(U)$.

\begin{definition}
	Let $R$ and $S$ be local rings with maximal ideals $\mm_R$ and $\mm_S$.
	A \vocab{morphism of local rings} is a homomorphism of rings
	$\psi : R \to S$ such that $\psi\pre(\mm_S) = \mm_R$.
\end{definition}
\begin{definition}
	A \vocab{morphism of locally ringed spaces}
	is a morphism of ringed spaces $f : (X, \OO_X) \to (Y, \OO_Y)$ such that
	for every point $p$ the induced map of stalks is a morphism of local rings.
\end{definition}
Recalling that $\mm_{X,p}$ is the maximal ideal of $\OO_X$ at $p$,
the new condition is saying that
\[ (f^\ast_P)\pre (\mm_{X,p}) = \mm_{Y,f(p)}. \]
Concretely, if $g$ is a function which vanishes on $f(p)$,
then its ``pullback'' $f^\ast_U(g)$ vanishes on $p$.

This completes the definition of a morphism of locally ringed spaces.
Isomorphisms of (locally) ringed spaces are defined in the usual way.

\section\problemhead
\begin{problem}
	Prove that if $\SF$ is a sheaf, then $\SF \cong \SF\sh$.
\end{problem}
\begin{sproblem}
	[Global section functor]
	We consider the category $\catname{ShRing}_X$ of sheaves of rings on a space $X$.
	The \vocab{global section functor}, denoted $\Gamma(X,-)$,
	which sends each sheaf $\SF$ to the ring $\SF(X)$.
	How does the global section functor behave on arrows?
	\begin{hint}
		Given a morphism of sheaves $\alpha : \SF \to \SG$
		read off a map $\SF(X) \to \SG(X)$.
	\end{hint}
\end{sproblem}

\chapter{Schemes}
Now that we understand sheaves well, we can readily define the scheme.
It will be a locally ringed space, so we need to define
\begin{itemize}
	\ii The set of points,
	\ii The topology on it, and
	\ii The structure sheaf on it.
\end{itemize}

In the case of $\Aff^n$, we used $\CC^N$ as the set of points
and $\CC[x_1, \dots, x_n]$ but then remarked that the set of points
of $\CC^n$ corresponded to the maximal ideals of $\CC[x_1, \dots, x_n]$.

In an \emph{affine scheme}, we will take an \emph{arbitrary} ring $R$,
and generate the entire structure from just $R$ itself.
The final result is called $\Spec R$, the \vocab{spectrum} of $R$.
The affine varieties $\VV(I)$ we met earlier will just be
$\CC[x_1, \dots, x_n] / I$, but now we will be able to take
\emph{any} ideal $I$, thus finally completing the table at the end
of the ``affine variety'' chapter.

In particular, $\Spec \CC[x] / (x^2)$ will be the double we sought for so long.

\section{The set of points}
\prototype{$\Spec \CC[x_1, \dots, x_n] / I$.}

First surprise, for a ring $R$:
\begin{moral}
	$\Spec R$ is defined as the set of prime ideals of $R$.
\end{moral}

This might be a little surprising, since we might have guessed
that $\Spec R$ should just have the maximal ideals.
What do the remaining ideals correspond to?
The answer is that they will be so-called \emph{generic points}
points which are ``somewhere'' in the space, but nowhere in particular.

\begin{remark}
	As usual $R$ itself is not a prime ideal, but $(0)$
	does if $R$ is an integral domain.
\end{remark}

\begin{example}
	[Examples of spectrums]
	\listhack
	\begin{enumerate}[(a)]
		\ii $\Spec \CC[x]$ consists of a point $(x-a)$ for every $a \in \CC$,
		which correspond to what we geometrically think of as $\Aff^1$.
		In additionally consists of a point $(0)$,
		which we think of as a ``generic point'', nowhere in particular.

		\ii $\Spec \CC[x,y]$ consists of points $(x-a,y-b)$
		(which are the maximal ideals) as well as $(0)$ again, a generic
		point that is thought of as ``somewhere in $\CC^2$,
		but nowhere in particular''.
		It also consists of generic points corresponding to irreducible
		polynomials $f(x,y)$, for example $(y-x^2)$,
		which is a ``generic point on the parabola''.

		\ii If $k$ is a field, $\Spec k$ is a single point,
		since the only maximal ideal of $k$ is $(0)$.
	\end{enumerate}
\end{example}
\begin{example}
	[Complex affine varieties]
	Let $I \subseteq \CC[x_1, \dots, x_n]$ be an ideal.
	Then \[ \Spec \CC[x_1, \dots, x_n] /I \] contains a
	point for every closed irreducible subvariety of $\VV(I)$.
	So in addition to the ``geometric points'' we have
	``generic points'' along each of the varieties.
\end{example}
\begin{example}
	[More examples of spectrums]
	\listhack
	\begin{enumerate}[(a)]
		\ii $\Spec \ZZ$ consists of a point for every prime $p$,
		plus a generic point that is somewhere, but no where in particular.

		\ii $\Spec \CC[x] / (x^2)$ has only $(x)$ as a prime ideal.
		The ideal $(0)$ is not prime since $0 = x \cdot x$.
		Thus as a \emph{topological space},
		$\Spec \CC[x] / (x^2)$ is a single point.
		
		\ii $\Spec \Zc{60}$ consists of three points.
		What are they?
	\end{enumerate}
\end{example}

\section{The Zariski topology of the spectrum}
\prototype{Still $\Spec \CC[x_1, \dots, x_n] / I$.}

Now, we endow a topology on $\Spec R$.
Since the points on $\Spec R$ are the prime ideals, we continue
the analogy by thinking of the points $f$ as functions on $\Spec R$. That is,
\begin{definition}
	Let $f \in R$ and $\pp \in \Spec R$.
	Then the \vocab{value} of $f$ at $\pp$ is defined to be $f \pmod{\pp}$.
	We denote it $f(\pp)$.
\end{definition}
\begin{example}
	[Vanishing locii in $\Aff^n$]
	Suppose $R = \CC[x_1, \dots, x_n]$,
	and $\pp = (x_1-a_1, x_2-a_2, \dots, x_n-a_n)$ is a maximal ideal of $R$.
	Then for a polynomial $f \in \CC$,
	\[ f \pmod \pp = f(a_1, \dots, a_n) \]
	with the identification that $\CC/\pp \cong \CC$.
\end{example}
Indeed if you replace $R$ with $\CC[x_1, \dots, x_n]$
and $\Spec R$ with $\Aff^n$ in everything that follows,
then everything will be clear.

\begin{definition}
	Let $f \in R$. We define the \vocab{vanishing locus} of $f$ to be
	\[ \VV(f) = \left\{ \pp \in \Spec R \mid f(\pp) = 0 \right\}
		= \left\{ \pp \in \Spec R \mid f \in \pp \right\}. \]
	More generally, just as in the affine case,
	we define the vanishing locus for an ideal $I$ as
	\begin{align*}
		\VV(I) &= \left\{ \pp \in \Spec R \mid f(\pp)=0 \forall f \in I \right\} \\
		&= \left\{ \pp \in \Spec R \mid f \in \pp \; \forall f \in I \right\} \\
		&= \left\{ p \in \Spec R \mid I \subseteq \pp \right\}.
	\end{align*}
	Finally, we define the \vocab{Zariski topology} on $\Spec R$
	by declaring that the sets of the form $\VV(I)$ are closed.
\end{definition}

Now, the following topological notion will come in handy:
\begin{definition}
	A point $p \in X$ is a \vocab{closed point} if the set $\{p\}$ is closed.
\end{definition}
\begin{ques}
	Show that $\pp \in \Spec R$ is a closed point
	if and only if $\pp$ is a maximal ideal.
\end{ques}
Therefore the Zariski topology lets us refer back to the old ``geometric''
as just the closed points.
\begin{example}
	[Generic points, continued]
	Let $R = \CC[x,y]$ and let $\pp = (y-x^2) \in \Spec \RR$;
	this is the ``generic point'' on a parabola.
	It is not closed, but we can compute its closure:
	\[
		\ol{\{\pp\}}
		= \VV(\pp) = \left\{ \qq \in \Spec R \mid \qq \supseteq \pp \right\}.
	\]
	This closure contains the point $\pp$ as well
	as several maximal ideals $\qq$, such as $(x-2,y-4)$ and $(x-3,y-9)$.
	In other words, the closure of the ``generic point'' of the parabola
	is literally the set of all points that are actually on the parabola
	(including generic points).

	That means the way to picture $\pp$ is a point that 
	is ``somewhere on the parabola'', but nowhere in particular.
	It makes sense then that if we take the closure,
	we get the entire parabola,
	since $\pp$ ``could have been'' any of those points.
\end{example}
The previous example illustrates the following important observation:
\begin{exercise}
	Let $I \subsetneq R$ be a proper ideal.
	Construct a bijection between the maximal ideals of $R$ containing $I$,
	and the maximal ideals of $R/I$.
	(\cite{ref:vakil} labels this exercise as
	``Essential Algebra Exercise (Mandatory if you haven't done it before)''.)
\end{exercise}

\begin{example}
	[The generic point of the $y$-axis isn't on the $x$-axis]
	Let $R = \CC[x,y]$ again.
	Consider $\VV(y)$, which is the $x$-axis of $\Spec R$.
	Then consider $\pp = (x)$, which is the generic point on the $y$-axis.
	Observe that
	\[ \pp \notin \VV(y). \]
	The geometric way of saying this is that a \emph{generic point}
	on the $y$-axis does not lie on the $x$-axis.
\end{example}

Finally, as before, we can define distinguished open sets,
which form a basis of the Zariski topology.
\begin{definition}
	Let $f \in \Spec R$.
	Then $D(f)$ is the set of $\pp$ such that $f(\pp) \neq 0$,
	a \vocab{distinguished open set}.
	These open sets form a basis for the Zariski topology on $\Spec R$.
\end{definition}

\section{The structure sheaf}
\prototype{Still $\CC[x_1, \dots, x_n] / I$.}

We have now endowed $\Spec R$ with the Zariski topology,
and so all that remains is to put a sheaf $\OO_{\Spec R}$ on it.
To do this we want a notion of ``regular functions'' as before.

In order for this to make sense, we have to talk about rational
functions in a meaningful way.
If $R$ was an integral domain, then we could just use the field of fractions;
for example if $R = \CC[x_1, \dots, x_n]$ then we could just
look at rational quotients of polynomials.

Unfortunately, in the general situation $R$ may not be an integral domain:
for example, the ring $R = \CC[x] / (x^2)$ corresponding to the double point.
So we will need to define something a little different:
we will construct the \emph{localization} of $R$ at a set $S$,
which we think of as the ``set of allowed denominators''.

\begin{definition}
	Let $S \subseteq R$, where $R$ is a ring,
	and assume $S$ is closed under multiplication.
	Then the \vocab{localization of $R$ at $S$}, denoted $S\inv R$,
	is defined as the set of fractions
	\[ \left\{ r/s \mid r \in R, s \in S \right\} \]
	where we declare two fractions $r_1 / s_1 = r_2 / s_2$ 
	to be equal if 
	\[ \exists s \in S : \quad s(r_1s_2 - r_2s_1) = 0. \]
\end{definition}
In particular, if $0 \in S$ then $S\inv R$ is the trivial ring.
So we usually only take situations where $0 \notin S$.
\begin{ques}
	Assume $R$ is an integral domain and $S = R \setminus \{0\}$.
	Show that $S\inv R$ is just the field of fractions.
\end{ques}

\begin{example}
	[Why the extra $s$?]
	The reason we need the condition $s(r_1s_2 - r_2s_1) = 0$
	rather than the simpler $r_1s_2 - r_2s_1 = 0$ is that
	otherwise the equivalence relation may fail to be transitive.
	Here is a counterexample: take
	\[ R = \Zc{12} \qquad S = \{ 0, 4, 8 \}. \]
	Then we have for example
	\[ \frac12 = \frac24 = \frac64 = \frac32. \]
	So we need to have $\frac12=\frac32$ which is only true
	with the first definition.

	Of course, if $R$ is an integral domain (and $0\notin S$)
	then this is a moot point.
\end{example}

The most important special case is the localization at a prime ideal.
\begin{definition}
	Let $R$ be a ring and $\pp$ a prime ideal.
	Then $R_\pp$ is defined to be $S\inv R$ for $S = R\setminus\pp$.
	We call this \vocab{localization at $\pp$}.
	Addition is defined in the obvious way.
\end{definition}
\begin{ques}
	Why is $S$ multiplicative closed in the above definition?
\end{ques}
Thus,
\begin{moral}
	If $R$ is functions on the space $\Spec R$,
	we think of $R_\pp$ as rational quotients $f/g$ where $g(\pp) \neq 0$.
\end{moral}
In particular, if $R = \CC[x_1, \dots, x_n]$
then this is precisely the definition of rational function from before!

Now, we can define the sheaf as ``locally rational'' functions.
This is done by a sheafification.
First, let $\SF$ be the pre-sheaf of ``globally rational'' functions:
i.e.\ we define $\SF(U)$ to be the following localization of $R$
to the functions vanishing outside $U$:
\[
	\SF(U) = \left\{ 
		\frac fg \mid f, g \in R
		\text{ and } g(\pp) \neq 0 \; \forall \pp \in U
	\right\}
	= \left(R \setminus \bigcup_{\pp \in U} \pp \right)\inv R.
\]
For every $\pp \in U$ we can view $f/g$ as an element in $R_\pp$
(since $g(\pp) \ne 0$).
As one might expect this is an isomorphism
\begin{lemma}[Stalks of the ``globally rational'' pre-sheaf]
	\label{lem:global_rational_stalk}
	The stalk of $\SF$ defined above $\SF_\pp$ is isomorphic to $R_\pp$.
\end{lemma}
\begin{proof}
	There is an obvious map $\SF_\pp \to R_\pp$ on germs by
	\[
		\left(U, f/g \in \SF(U) \right)
		\mapsto f/g \in R_\pp . \]
	(Note the $f/g$ on the left lives in
	$\SF(U) = \left(R \setminus \bigcup_{\pp \in U} \pp \right)\inv R$
	but the one on the right lives in $R_\pp$).
	Now suppose $(U_1, f_1 / g_1)$ and $(U_2, f_2 / g_2)$
	are germs with $f_1/g_1 = f_2/g_2 \in R_\pp$.
	\begin{exercise}
		Now, show both germs are equal to $(U_1 \cap U_2 \cap D(h), f_1 / g_1)$
		with $D(h)$ the distinguished open set.
	\end{exercise}
	It is also surjective, since given $f/g \in R_\pp$ we take $U = D(g)$
	the distinguished open set for $g$.
\end{proof}

Then, we set \[ \OO_{\Spec R} = \SF\sh. \]
In fact, we can even write out the definition of the sheafification,
by viewing the germ at each point as an element of $R_\pp$.
\begin{definition}
	Let $R$ be a ring. Then $\Spec R$ is made into a ringed space by setting
	\[ \OO_{\Spec R}(U) 
		= \left\{ (f_\pp \in R_\pp)_{\pp \in U}
		\text{ which are locally $f/g$} \right\}. \]
	That is, it consists of sequence $(f_\pp)_{\pp \in U}$, with
	each $f_\pp \in R_\pp$, such that for every point $\pp$ there
	is a neighborhood $U_\pp$ and an $f,g \in R$ such that
	$f_\qq = \frac fg \in R_\qq$ for all $\qq \in U_\pp$.
\end{definition}

\section{Properties of affine schemes}
Now that we're done defining an affine scheme,
we state some important results about them.

\begin{definition}
	For $g \in R$, we define the \vocab{localization of $R$ at $g$},
	denoted $R_g$ to be $\{1, g, g^2, g^3, \dots\}\inv R$.
	(Note that $\left\{ 1, g, g^2, \dots \right\}$ is multiplicatively closed.)
\end{definition}
This is admittedly somewhat terrible notation, since $R_\pp$
is the localization with multiplicative set $R \setminus \pp$,
while $R_g$ is the localization with multiplicative set $\{1,g,g^2,\dots\}$;
these two are quite different beasts!

\begin{example}
	[Localization at an element]
	Let $R = \CC[x,y,z]$ and let $g = x$.
	Then
	\[ R_g = \left\{ \frac{P(x,y,z)}{x^n} \mid
		P \in \CC[x,y,z], \; n \ge 0 \right\}. \]
\end{example}
\begin{theorem}
	[On the affine structure sheaf]
	Let $R$ be a ring and $\Spec R$ the associated affine scheme.
	\begin{enumerate}[(a)]
		\ii Let $\pp$ be a prime ideal.
		Then $\OO_{\Spec R, \pp} \cong R_\pp$.
		\ii Let $D(g)$ be a distinguished open set.
		Then $\OO_{\Spec S}(D(g)) \cong R_g$.
	\end{enumerate}
\end{theorem}
This matches the results that we've seen when $\Spec R$ is an affine variety.
\begin{proof}
	Part (a) follows by \Cref{lem:global_rational_stalk}
	and \Cref{lem:pre_sheaf_stalk}.
	For part (b), we need the following:
	\begin{ques}
		If $I$ and $J$ are ideals of a ring $R$,
		then $\VV(I) \subseteq \VV(J)$ if and only if
		$\sqrt{J} \subseteq \sqrt{I}$.
		(Use the fact that $\sqrt I = \cap_{I \subseteq \pp} \pp$.)
	\end{ques}
	Then we can repeat the proof of \Cref{thm:reg_func_distinguish_open}.
\end{proof}

\begin{theorem}
	[Affine schemes and commutative rings are the same category]
	Let $R$ and $S$ be rings.
	There is a natural bijection between maps of schemes $\Spec R \to \Spec S$	
	and ring homomorphisms $\psi : S \to R$.
\end{theorem}
\begin{proof}
	First, we need to do the construction:
	given a map of schemes $f : \Spec R \to \Spec S$ 
	we need to construct a homomorphism $S \to R$.
	This should be the easy part, because $f$ has a lot of data.
	Indeed, recall that for every $U \subseteq \Spec S$
	there is supposed to be a map
	\[ f^\ast_U : \OO_{\Spec S} (U) \to \OO_{\Spec R}(f\pre(U)). \]
	If we take $U$ to be the entire space $\Spec S$ we now have a ring homomorphism
	\[ S = \OO_{\Spec S} (\Spec S) \to \OO_{\Spec R} (\Spec R) = R. \]
	which is one part of the construction.

	In fact, this is just the global section functor
	applied to $\OO_{\Spec S} \to f_\ast \OO_{\Spec R}$.

	The more involved direction is building from a homomorphism
	$\psi : S \to R$ a map $f : \Spec R \to \Spec S$.
	We define it by:
	\begin{itemize}
		\ii On points, $f(\pp) = \psi\pre(\pp)$ for each prime ideal $\pp$.
		One can check $\psi\pre(\pp)$ is indeed prime
		and that $f$ is continuous in the Zariski topology.
		\ii We first construct a map on the stalks of the sheaf as follows:
		for every prime ideal $\qq$ in $R$, let $\psi``(\qq) = \pp$ be its image
		(so that $f(\qq) = \pp$) and consider the map
		\[ \psi_\qq : \OO_{\Spec S, \qq} = S_{\qq}
			\to R_{\psi``(\qq)} = \OO_{\Spec R, \psi``(\qq)}
			\quad\text{by}\quad f/g \mapsto \psi(f)/\psi(g). \]
		Again, one can check this is well-defined and a map of local rings.
		This is called the localization of $\psi$ at $\qq$.
		\ii Finally, we use the above map to construct an
		$f^\ast_U : \OO_{\Spec S}(U) \to \OO_{\Spec R}(f\pre(U))$
		for every open set  $U \subseteq \Spec S$.
		Since $\OO_{\Spec S}$ is a sheafification, we think of its
		sections as sequences of compatible germs $(g_\qq)_{\qq \in U}$
		and then map it via the map $\psi_\qq$ above.
	\end{itemize}
	One then has to check that everything is well-defined.
	This is left as an exercise to the overzealous reader;
	for an actual proof see \cite[Proposition 6.3.2]{ref:vakil}.
\end{proof}
\begin{remark}
	Categorically, this says that the ``global section functor''
	on the category $\catname{AffSch}$ of affine schemes is a
	fully faithful functor to $\catname{CRing}$.
	Moreover, it is surjective on objects
	(the correct term here is ``essentially surjective on objects'').
\end{remark}
To be more philosophical,
\begin{moral}
	The category of affine schemes and the
	category of commutative rings are \emph{exactly the same},
	down to the morphisms between two pairs of objects.
\end{moral}


\section{Schemes}
Now we can define a general scheme:
\begin{definition}
	A \vocab{scheme} is a locally ringed space $(X, \OO_X)$
	with an open affine cover $\{U_\alpha\}$ of $X$
	such that each pair $(U_\alpha, \OO_X \restrict{U_\alpha})$
	is isomorphic to an affine scheme.
\end{definition}
Hooray!

\section{Projective scheme}
\prototype{Projective varieties, in the same way.}
The most important class of schemes which are not affine are
\emph{projective} schemes.
The complete the obvious analogy:
\[
	\frac{\text{Affine variety}}{\text{Projective variety}}
	= 
	\frac{\text{Affine scheme}}{\text{Projective scheme}}.
\]
Let $S$ be \emph{any} (commutative) graded ring, like $\CC[x_0, \dots, x_n]$.
\begin{definition}
	We define $\Proj S$, the \vocab{projective scheme over $S$}:
	\begin{itemize}
		\ii As a set, $\Proj S$ consists of \emph{homogeneous prime ideals}
		$\pp$ which do not contain $S^+$.
		\ii If $I \subseteq S$ is homogeneous, then
		we let $\Vp(I) = \{ \pp \in \Proj S \mid \pp \subseteq I \}$.
		Then the \vocab{Zariski topology} is imposed by declaring 
		sets of the form $\Vp(I)$ to be closed.
		\ii We now define a pre-sheaf $\SF$ on $\Proj S$ by
		\[ \SF(U) = 
			\left\{ \frac{f}{g} \mid 
			g(\pp) \neq 0 \; \forall \pp \in U \text{ and }
			\deg f = \deg g \right\}.
		\]
		In other words, the rational functions are quotients $f/g$
		where $f$ and $g$ are \emph{homogeneous of the same degree}.
		Then we let \[ \OO_{\Proj S} = \SF\sh \] be the sheafification.
	\end{itemize}
\end{definition}
\begin{definition}
	The \vocab{distinguished open sets} $D(f)$ of the $\Proj S$
	are defined as $\left\{ \pp \in \Proj S : f(\pp) \neq 0 \right\}$,
	as before; these form a basis for the Zariski topology of $\Proj S$.
\end{definition}
Now, we want analogous results as we did for affine structure sheaf.
So, we define a slightly modified localization:
\begin{definition}
	Let $S$ be a graded ring.
	\begin{enumerate}[(i)]
		\ii For a prime ideal $\pp$, let
		\[ S_{(\pp)} = \left\{ \frac fg \mid g(\pp) \neq 0 \text{ and }
			\deg f = \deg g \right\} \]
		denote the elements of $S_\pp$ with ``degree zero''.
		\ii For any homogeneous $g \in S$ of degree $d$, let
		\[ S_{(g)} = \left\{ \frac{f}{g^r} \mid 
			\deg f = r \deg g \right\} \]
		denote the elements of $S_g$ with ``degree zero''.
	\end{enumerate}
\end{definition}

\begin{theorem}
	[On the projective structure sheaf]
	Let $S$ be a graded ring and let $\Proj S$ the associated projective scheme.
	\begin{enumerate}[(a)]
		\ii Let $\pp \in \Proj S$.
		Then $\OO_{\Proj S, \pp} \cong S_{(\pp)}$.
		\ii Suppose $g$ is homogeneous with $\deg g > 0$. Then
		\[ D(g) \cong \Spec S_{(g)} \]
		as locally ringed spaces.
		In particular, $\OO_{\Proj S}(D(g)) \cong S_{(g)}$.
	\end{enumerate}
\end{theorem}
\begin{ques}
	Conclude that $\Proj S$ is a scheme.
\end{ques}

Of course, the archetypal example is that
\[ \Proj \CC[x_0, x_1, \dots, x_n] / I \]
corresponds to the projective subvariety of $\CP^n$
cut out by $I$ (when $I$ is radical).
In the general case of an arbitrary ideal $I$, we
call such schemes \vocab{projective subscheme} of $\CP^n$
For example, the ``double point'' is given by $\Proj[x_0,x_1]/(x_0^2)$.

\begin{remark}
	No comment yet on what the global sections of $\OO_{\Proj S}(\Proj S)$ are.
	(The theorem above requires $\deg g > 0$, so we cannot just take $g=1$.)
	One might hope that in general $\OO_{\Proj S}(\Proj S) \cong S^0$
	in analogy to our complex projective varieties, but
	one needs some additional assumptions on $S$ for this to hold.
\end{remark}

\section{Fun application: Chinese remainder theorem}
\todo{write me}

\section\problemhead

\chapter{More on sheaves}
Previously, we defined a sheaf so that we could equip
it onto a scheme and thus produce a ringed space.
However, we will now proceed to define some additional sheaves on a scheme $X$,
and study the relations between them.

\section{Sheaves of $\OO_X$-modules}
\prototype{Twisting Sheaves}
Up until now we have mostly been working with sheaves of rings,
but now we want to instead think about sheaves of modules related
to a scheme $X$. 
\begin{definition}
	Let $(X, \OO_X)$ be a ringed space.
	A \vocab{sheaf of $\OO_X$-modules} is a sheaf $\SF$
	of abelian groups such that $\SF(U)$ is also an $\OO_X(U)$-module.
	The module structures are also required to be compatible
	with the restriction maps.
\end{definition}
\begin{example}
	[Twisting sheaves]
	Let $X = \Proj S$ be a projective scheme for a graded ring $S$
	(e.g.\ a projective subscheme).
	For any integer $m$ we consider a pre-sheaf
	\[ \SF^n = \left\{ \frac fg \mid g(\pp)\neq 0 \;\forall p \in X 
		\text{ and } \deg f - \deg g = n \right\}.
	\]
	Intuitively, these are degree $n$ ``functions''.
	Then $\OO_X(n) = (\SF^n)\sh$;
	we call these the \vocab{twisting sheaves}.

	For example, if $S = \CC[x_0, \dots, x_n]$ then
	\begin{enumerate}[(a)]
		\ii $\OO_X(0) \cong \OO_X$, by definition.
		\ii A degree $n$ polynomial defines a global section of $\OO_X(n)$.
		\ii $\OO_X(m)$ has no nonzero global sections if $m < 0$.
		\ii Suppose $X = \CP^1$ has coordinates $(s:t)$. Then
		\[ \frac 1s \in \OO_X(-1)(U_0) \]
		is an example of a (non-global) section of $\OO_X(-3)$.
		\ii Suppose $X = \CP^2$ has coordinates $(x:y:z)$. Then
		\[ \frac{1}{xyz} \in \OO_X(-1)(U_0 \cap U_1 \cap U_2) \]
		is an example of a (non-global) section of $\OO_X(-3)$.
	\end{enumerate}
\end{example}

Observe that the stalks of such a sheaf are abelian groups.

\section{Sheaves of $\OO_X$-modules make an abelian category}
I won't take the time to do this properly,
since I won't really use much of it.
See \cite{ref:vakil} for a proper exposition.

Fix a ringed space $(X, \OO_X)$.
It's easy to see there is a category of sheaves of $\OO_X$-modules.
(the objects are sheaves and the morphisms are, well, morphisms of sheaves).
To see that it is abelian, we need to do the following:
\begin{itemize}
	\ii Identify the zero object of the category.
	This is constant sheaf $0$, which gives $\SF(U) = \{0\}$ for every $U$.
	(This more or less follows from the fact that $\{0\}$
	is the zero object of the abelian category $\catname{AbGrp}$.)
	\ii Show that we can add two morphisms.
	This is done in the obvious manner.
	\ii Show that any morphism has a ``kernel sheaf'' and ``cokernel sheaf''.
	The definition is not too bad: given $\alpha : \SF \to \SG$ we define the
	pre-sheaves $\ker\alpha$ and $\coker'\alpha$ by
	\begin{align*}
		\ker(\alpha)(U) &= \ker(\alpha_U) \\
		\coker'(\alpha)(U) &= \coker(\alpha_U).
	\end{align*}
	It turns out that $\ker(\alpha)$ is already a sheaf.
	But $\coker'(\alpha)$ is only a pre-sheaf and in general not a sheaf.
	(Again, see \cite{ref:vakil} or \cite{ref:gathmann} for why.)
	Thus we have to define $\coker(\alpha) = \coker'(\alpha) \sh$.
	
	The tricky part, done in \cite{ref:vakil} and the one which I'm omitting,
	is showing that this kernel and cokernel sheaf satisfy
	the correct universal properties.
\end{itemize}

From this we deduce
\begin{theorem}
	[Sheaves form an abelian category]
	Let $(X, \OO_X)$ be a ringed space.
	The category of sheaves of $\OO_X$-modules forms an abelian category.
\end{theorem}

In particular, we can talk about whether a sequence of sheaves
$\SF \to \SG \to \SH$ is exact: just use the categorical definition.
However, if possible we want to avoid using the long-winded definition.
So, I will just quote the relevant results:
\begin{theorem}
	[Exactness can be checked at stalks]
	A sequence of sheaves of $\OO_X$-modules
	$\SF \to \SG \to \SH$ is exact if and only if
	for every point $p \in X$,
	the induced map of stalks $\SF_p \to \SG_p \to \SH_p$ is exact.
\end{theorem}
Thus, we \emph{really} like stalks; they behave as nicely as possible.
\begin{ques}
	Deduce that $\SF \to \SG$ is injective / surjective / an isomorphism
	if and only if $\SF_p \to \SG_p$ is injective / surjective / an isomorphism
	for every point $p$.
\end{ques}

If we insist on sections, the only useful result we have is:
\begin{theorem}
	[Monic/epic on sections]
	Let $\alpha : \SF \to \SG$ be a morphism of $\OO_X$-modules.
	\begin{enumerate}[(a)]
		\ii The map $\alpha$ is monic if and only if $\SF(U) \to \SG(U)$
		is injective for every open $U$.
		\ii The map $\alpha$ is epic if $\SF(U) \to \SG(U)$
		is surjective for every open $U$,
		but the converse is not true in general.
	\end{enumerate}
\end{theorem}

In addition, we can do the following constructions.
\begin{definition}
	If $\SF$ and $\SG$ are sheaves of $\OO_X$-modules, then
	\begin{enumerate}[(a)]
		\ii The \vocab{direct sum} $\SF \oplus \SG$ is defined
		on open sets by \[ U \mapsto \SF(U) \oplus \SG(U). \]
		(This is already a sheaf, so we don't need to sheafify.)
		\ii The \vocab{tensor product} $\SF \otimes \SG$ is defined
		by taking the sheafification of the pre-sheaf
		\[ U \mapsto \SF(U) \otimes_{\OO_X(U)} \SG(U). \]
		\ii The \vocab{dual} $\SF^\vee$ of $\SF$ is defined
		by taking the sheafification of the pre-sheaf
		\[ U \mapsto \Hom_{\OO_X(U) \text{ module}} (\SF(U), \OO_X(U)). \]
	\end{enumerate}
\end{definition}
\begin{example}[More on twisting sheaves]
	Consider the twisting sheaves of $X = \CP^N$.
	\begin{enumerate}[(a)]
		\ii $\OO_X(n) \otimes \OO_X(m) \cong \OO_X(m+n)$.
		What do you think the isomorphism is?
		\ii $\OO_X(n)^\vee \cong \OO_X(-n)$.
	\end{enumerate}
\end{example}

\section{Quasi-coherent sheaves}

\chapter{Divisors and line bundles}

\end{document}
